\chapter{The Supraverse}
\label{chap:supraverse}

\section{Global structure}

The supraverse is a single 4D manifold (3 space $+$ 1 time, Lorentzian
signature) with non-trivial topology. Most of it is essentially vacuum---
quantum-mechanically active but unstructured. Embedded in this vacuum are
condensed regions: universes, each a connected sub-manifold with its own
internal cosmological geometry, joined to the rest by Cartasis hypersurfaces.

The topology is foam-like: many universe-bubbles distributed across the
supraverse, connected to each other through nested parent--child relationships
via bounce membranes.

\section{No fifth dimension required}

The supraverse is genuinely 4D. Earlier intuitions about needing an extra
dimension to ``embed'' the parent--child relationship are visualization
crutches, not physical requirements. Mathematically, the topology is well-defined
intrinsically: two universes connected at a Cartasis hypersurface is a 4D
manifold with a particular boundary identification, no ambient space needed.

For visualization purposes, one might embed the supraverse in higher dimensions
(the Whitney embedding theorem guarantees this is possible). But that embedding
has no physical meaning---it is a drawing aid.

\section{Equilibrium configuration}

\subsection{Why condensation is favored}

Naively, one might expect universe formation to be entropically disfavored---
universes are ``ordered'' structures emerging from ``disordered'' vacuum. This
intuition is wrong because it ignores gravity. Gravitational entropy has
inverted relations compared to ordinary thermodynamics:

\begin{itemize}
  \item uniform matter distribution: low gravitational entropy;
  \item clumped matter: higher gravitational entropy;
  \item black holes: maximum entropy for given mass--energy
    (Bekenstein--Hawking, $S \propto A/4\ell_P^2$)~\citep{bekenstein1973}.
\end{itemize}

Penrose has emphasized this for decades~\citep{penrose1979}. Applied to the
supraverse: uniform vacuum has zero gravitational structure and low
gravitational entropy; a condensed foam of universes---each containing black
holes, structure, horizons---has high gravitational entropy.

\textbf{The Second Law drives the supraverse toward condensation, not away from
it.} Spontaneous universe formation is entropy-increasing when gravity is
properly accounted for. The supraverse must be in a foam state at equilibrium.

\subsection{What the equilibrium looks like}

The supraverse in its equilibrium configuration contains:

\begin{itemize}
  \item a vast quantum-vacuum background, essentially everywhere;
  \item sparse, randomly distributed universe-bubbles within this background;
  \item each universe a bubble containing its own internal cosmological structure;
  \item universes connected to each other through nested Cartasis membranes
    (parent--child relationships);
  \item a distribution of universe sizes peaking at a thermodynamically optimal
    value (see below);
  \item matter-biased and antimatter-biased lineages in equal numbers (CPT
    symmetry preserved globally).
\end{itemize}

This is the stationary distribution. Local universes are born, evolve, produce
descendants, and eventually evaporate via parent Hawking processes. New OG
universes nucleate from vacuum fluctuations. The supraverse maintains its
statistical distribution across all this dynamics.

\section{Universe size distribution}

\subsection{The competing factors}

The production rate of OG universes of mass $M$ depends on two factors. The
\emph{fluctuation probability}---the probability of a vacuum fluctuation
producing Cartasis density over a region containing mass $M$---scales roughly as
$\exp(-M c^2 \tau / \hbar)$ for some characteristic timescale $\tau$: smaller is
more probable. The \emph{descendant productivity}---a baby universe of mass $M$
has Hawking lifetime $\tau_H \propto M^3$, during which it can produce many
descendant black holes---means larger $M$ is more productive.

\subsection{The optimum}

The peak of the supraverse population distribution is where these factors
balance. The exact location depends on Einstein--Cartan parameters, but
order-of-magnitude estimates place it somewhere in the range of cosmological
universes---large enough for substantial astrophysics, small enough that the
fluctuation probability has not completely tanked.

\begin{conjecture}
The peak is near our observed universe's mass. Universes much smaller fail to
produce many descendants; universes much larger are exponentially rare. We sit
near the optimum.
\end{conjecture}

This is a quantitative prediction. Simulations of supraverse dynamics could in
principle compute the peak location and check whether our parameters are near it
(Chapter~\ref{chap:sims}).

\section{Lineages and the family tree}

Each OG universe seeds a lineage tree of descendants. The tree branches at each
black-hole formation within a universe (each black hole potentially produces a
baby universe).

\subsection{Branching structure}

A typical viable universe contains many stellar-mass black holes ($\sim 10^{18}$
in a galaxy-rich universe like ours), some supermassive black holes
($\sim 10^{9}$ across all galaxies), and possibly primordial black holes from
early-universe dynamics. Each is a potential baby universe, so a typical lineage
tree branches widely at each generation. The tree depth is limited by the finite
Hawking lifetime of each generation, by the requirement that each generation be
viable enough to produce its own black holes, and by the cumulative effects of
flip $+$ filter on chirality.

\subsection{Lineage chirality structure}

Each universe acquires its matter--antimatter bias \emph{at its own bounce}, not
by inheritance down a lineage. The Einstein--Cartan bounce is CPT-even
(Chapter~\ref{chap:chirality}): it neither flips charge nor alternates chirality
across generations. Instead, the parity-violating chiral-vortical transport at
each rotating bounce freezes in a bias $|\eta|\sim C\,|\omega/T|$ whose
\emph{sign} is set by that bounce's vorticity. Because vorticity sign is random
across the supraverse, matter- and antimatter-biased universes occur in equal
numbers and global CPT is preserved. There is no alternating chirality to climb
and no compounding filter; the bias is made fresh at every membrane crossing.

\textbf{Observer placement:} observers exist in universes whose birth bias clears
the viability threshold---enough asymmetry to leave matter after annihilation and
form structure. This is a one-step \emph{birth filter} on bounce vorticity
($|\omega/T| > \eta_{\min}/C$), not a position in an alternating lineage.

\section{OG bounce nucleation rate}

The rate of OG bounce nucleation depends on the vacuum-fluctuation spectrum at
high energies (set by QFT), the Cartasis density threshold (set by
Einstein--Cartan parameters), and the available supraverse spacetime volume.

For our supraverse to contain even one universe like ours, the cumulative
nucleation rate over supraverse history must be at least one per ``us-sized''
mass scale. For the supraverse to contain many universes, the rate must be
higher. The exact rate is calculable in principle from QFT vacuum-fluctuation
statistics plus Einstein--Cartan thresholds. It is likely extremely small per
unit spacetime volume, but the supraverse compensates by being effectively
infinite.

\section{Genesis, growth, and the absence of hell}

\subsection{No hell upward}

An older worry held that, climbing the lineage toward earlier ancestors, a
compounding filter would make ancestors progressively more symmetric---less
asymmetry, more annihilation, until near-symmetric, sterile universes that cannot
form black holes (the ``hell upward'' problem). With the CPT-even bounce this
dissolves on two counts. First, there is no charge flip and so no alternating
chirality to climb (Chapter~\ref{chap:chirality}). Second, and decisively, the
asymmetry is \emph{generated at every bounce}, not diluted along a lineage: each
rotating crossing writes $|\eta|\sim C|\omega/T|$ afresh. A universe's bias is set
by its own bounce, not by its distance from some symmetric ancestor. The upward
end of the chain is therefore not brimstone---it is the void
(Section~\ref{sec:genesis-void}), and the OG universes born from it are
immediately viable provided their bounce rotated enough.

\subsection{The void and OG genesis}
\label{sec:genesis-void}

An OG universe is nucleated directly from the primordial void, with no parent. It
cannot inherit a bias, so its asymmetry must be \emph{made} at the OG bounce
itself---and the rotating-bounce chiral-vortical filter does exactly that,
writing matter at the very first crossing. The void is the
low-gravitational-entropy, CPT-symmetric substrate; genesis is the Second Law
acting on it, condensing rotating fluctuations into bounces. Most fluctuations
are too small or too slowly rotating to clear the birth filter and produce only
sterile radiation; the rare viable ones seed everything. \emph{Most supraverse
structure is concentrated in the successful OG lineages}; failed attempts do not
propagate.

\subsection{Growth: biased runaway accretion}

The condensation favoured by the Second Law has a concrete microphysics: black
holes have negative heat capacity, so above a critical mass
$M_{\mathrm{crit}} = \hbar c^3/(8\pi G k_B T_{\mathrm{bath}})$ a hole accretes
faster than it Hawking-radiates and grows. (In today's CMB bath
$M_{\mathrm{crit}}\sim 5\times10^{22}\,\mathrm{kg}$, so every astrophysical black
hole is already growing.) Because the Cartasis interface is persistent
(Chapter~\ref{chap:baby}), the parity-violating filter keeps acting throughout
the parent's life, so the growth is \emph{biased}---the hole preferentially
retains one charge sign of what it accretes. Two consequences: baryogenesis is
\emph{ongoing} rather than one-shot (helping the single-bounce shortfall of
Chapter~\ref{chap:chirality}), and the migration of the void's free energy into
the growing universe-holes \emph{is} the condensation that makes the foam its
equilibrium. From inside a baby universe this continuous biased injection is the
dark sector (Chapter~\ref{chap:darksector}).

\subsection{Infinitely many OGs; do they grow forever?}

By Axiom~\ref{ax:qm} plus infinity, the supraverse nucleates \emph{infinitely
many} OG universes. Whether an individual OG grows without bound is open. In a
static infinite bath the negative heat capacity drives unbounded runaway
accretion; in practice growth is regulated---by depletion of the local
environment, by the void's own expansion diluting the bath, and by the stationary
balance of the foam (large holes growing, new OGs nucleating, sub-critical holes
evaporating). The honest statement is that individual universes can grow very
large while the population remains statistically stationary; the global question
is entangled with the measure problem (Appendix~\ref{app:openq}).

\section{Information structure of the supraverse}

\subsection{Local vs.\ global unitarity}

Each universe, viewed in isolation, appears to lose information through Hawking
radiation. But the Hawking radiation is entangled with the universe's interior
content, and that entanglement extends across bounce membranes.

\textbf{The supraverse as a whole is unitary.} Information that ``disappears''
from a baby universe via Hawking radiation reappears in the parent's exterior.
Information that flows into a parent black hole reappears in the baby's interior.
Total information is conserved across the supraverse. This resolves the
black-hole information paradox at the supraverse level: each universe's internal
unitarity may be approximate, but global supraverse unitarity is exact.

\subsection{Entanglement across membranes}

The persistent Cartasis channel maintains entanglement between parent exterior
and baby interior throughout the parent's lifetime. This is the structural
realization of ER\,$=$\,EPR~\citep{maldacena2013}: the geometric connection
between parent and baby is the entanglement structure between Hawking radiation
and the baby's matter content.

\section{Time within the supraverse}

There is no privileged supraverse-wide time direction. Each universe has its own
internal time arrow (set by its bounce-to-future thermodynamic gradient). These
arrows point in different directions in supraverse coordinates. The supraverse as
a whole is globally time-symmetric: equal numbers of universes with arrows
pointing ``this way'' and ``that way.'' See Chapter~\ref{chap:timearrow} for the
detailed treatment.
